\subsection{Overview}

Lossless image compression has been a fundamental area of research in signal processing and information theory for decades. The primary goal is to reduce the number of bits required to represent an image without any loss of information, enabling perfect reconstruction of the original image.

\subsection{Entropy Coding Methods}

\textbf{Huffman coding} \cite{huffman1952} assigns variable‑length codes to symbols based on their frequencies, achieving optimal coding for a given probability distribution when using integer‑length codes. \textbf{Arithmetic coding} \cite{witten1987arithmetic} extends this by allowing fractional bit allocations, often yielding better compression than Huffman coding, albeit with higher computational complexity. \textbf{Lempel–Ziv–Welch (LZW)} \cite{welch1984technique} is a dictionary‑based method that replaces repeated strings with references to a dictionary built adaptively during compression.

\subsection{Predictive and Transform Coding}

\textbf{JPEG‑LS} (ISO/IEC 14495‑1) is an international standard for lossless compression of continuous‑tone images. It employs a median edge detector (MED) predictor and Golomb–Rice coding \cite{weinberger2000loco}. \textbf{CALIC} (Context‑based Adaptive Lossless Image Coding) \cite{wu1996context} uses gradient‑based prediction and context modeling of prediction errors, achieving state‑of‑the‑art compression ratios.

\subsection{Transform Domain Methods}

The \textbf{S‑Transform} and \textbf{S+P Transform} \cite{said1996new} provide reversible integer‑to‑integer transforms. The S+P transform combines a Haar‑like decomposition with a prediction step to reduce entropy. \textbf{Integer wavelet transforms} (IWTs) using the lifting scheme \cite{calderbank1998wavelet} enable perfect reconstruction with integer arithmetic.

\subsection{Clustering and Neural Network Approaches}

\textbf{K‑Means clustering} and \textbf{Gath–Geva fuzzy clustering} have been applied to color quantization. Recent works have explored \textbf{Fuzzy Backpropagation Neural Networks (FBPNN)} that integrate fuzzy membership values into weight updates \cite{pedrycz1996fuzzy}.